\section{Appendix to Section 5: Monotonicity and Threshold Existence}

This appendix provides the technical details behind the sign conditions and
threshold objects used in Section~5. Throughout, all expressions are based on
the exact stage-4 first-order conditions reported in Section~3 and on the
continuation-welfare formulas collected in Appendix~A.

\subsection{Derivative objects and benchmark signs}

Recall the continuation-welfare differences
\begin{align}
\Delta_A^I &= \widehat W_A^{IS}-\widehat W_A^{SU}, &
\Delta_C^I &= \widehat W_C^{IS}-\widehat W_C^{SU}, \label{eq:appendix_deltas_1}\\
\Delta_A^{SU} &= \widehat W_A^{SU}-\widehat W_A^{SW}, &
\Delta_A^{IS} &= \widehat W_A^{IS}-\widehat W_A^{SW}. \label{eq:appendix_deltas_2}
\end{align}

The derivatives used in Section~5 are obtained by differentiating the exact
continuation-welfare formulas in Appendix~A. The resulting rational functions
are lengthy, so the fully expanded numerators are omitted here. What matters
for Section~5 is the sign of those derivatives on the compact parameter domain
used for the policy discussion.

A useful benchmark check is the case without network effects. Evaluating the
derivatives at $v=0$ gives
\begin{align}
\left.\frac{\partial \Delta_A^I}{\partial \tau_C}\right|_{v=0}
&= -\frac{5c}{16}, \label{eq:dAI_tauC_v0_final}\\
\left.\frac{\partial \Delta_C^I}{\partial \tau_C}\right|_{v=0}
&= -\frac{9c}{4}, \label{eq:dCI_tauC_v0_final}\\
\left.\frac{\partial \Delta_A^{SU}}{\partial \tau}\right|_{v=0}
&= -\frac{21c}{16}, \label{eq:dASU_tau_v0_final}\\
\left.\frac{\partial \Delta_A^{IS}}{\partial \tau}\right|_{v=0}
&= -\frac{13c}{8}. \label{eq:dAIS_tau_v0_final}
\end{align}
These expressions are strictly negative whenever $c>0$.

For later comparative statics, it is also useful to note that
\begin{equation}
\left.\frac{\partial \Delta_A^I}{\partial \tau}\right|_{v=0}
= -\frac{5c}{16}<0.
\label{eq:dAI_tau_v0_final}
\end{equation}

\subsection{A broader numerically verified monotonicity domain}

Because the exact derivative formulas are high-order rational functions, the
most transparent way to support the policy analysis in Section~5 is to verify
their signs numerically on an explicit compact parameter domain.

We use the broader compact domain
\begin{equation}
\mathcal D'
=
\left\{
\begin{aligned}
(v,c,\tau,\tau_C):\quad &0 \le v \le 0.055,\\
&0.03 \le c \le 0.06,\\
&0 \le \tau \le 0.04,\\
&0 \le \tau_C \le 0.18
\end{aligned}
\right\}
\subset \mathcal P.
\label{eq:domain_D}
\end{equation}
This domain contains the benchmark examples used in Section~4 and in the
policy discussion of Section~5, while allowing a materially broader range of
network-effect and baseline trade-cost values than the narrower benchmark box
used to motivate the numerical exercises.

On a $31\times 31\times 31\times 31$ uniform grid over $\mathcal D'$, all
admissibility conditions defining $\mathcal P$ remain strictly satisfied. In
particular,
\begin{align}
\min_{\mathcal D'}\left(\frac{2}{7}-v\right) &= 0.2307142857, \notag\\
\min_{\mathcal D'}\left(1-2c-3\tau+\tau_C\right) &= 0.70, \notag\\
\min_{\mathcal D'}\left(1-2c+\tau-3\tau_C\right) &= 0.34, \notag\\
\min_{\mathcal D'}\left(1-2c-2\tau\right) &= 0.80, \label{eq:D_admissibility_margins}
\end{align}
and the smallest regime-specific positivity margin on the grid is
\begin{equation}
\min_{\mathcal D'}\Xi_z^{SU}=0.015093>0.
\label{eq:D_smallest_margin}
\end{equation}
Hence the entire verified domain lies strictly inside the admissible set.

The same grid check yields
\begin{align}
\max_{\mathcal D'}\frac{\partial \Delta_A^I}{\partial \tau_C}
&= -0.009375, \label{eq:D_dAI_max}\\
\max_{\mathcal D'}\frac{\partial \Delta_C^I}{\partial \tau_C}
&= -0.0675, \label{eq:D_dCI_max}\\
\max_{\mathcal D'}\frac{\partial \Delta_A^{SU}}{\partial \tau}
&= -0.039375, \label{eq:D_dASU_max}\\
\max_{\mathcal D'}\frac{\partial \Delta_A^{IS}}{\partial \tau}
&= -0.04875. \label{eq:D_dAIS_max}
\end{align}
Since even the maxima are strictly negative, the monotonicity conditions used
in Section~5 hold throughout the grid. The same sign pattern was also
confirmed by an additional random check with 300{,}000 draws from
$\mathcal D'$.

Accordingly, for the parameter region used in the paper's numerical policy
discussion, the threshold analysis in Section~5 is conducted on a compact
domain on which the required monotonicity is numerically verified rather than
assumed globally over all of $\mathcal P$. The policy results should still be
read as local-to-domain comparative statics rather than as unrestricted global
theorems, but the verified region is not confined to an arbitrarily thin
neighborhood of the benchmark point.

\subsection{Threshold existence: generic lemma}

The threshold objects in Section~5 are applications of the following
elementary result.

\begin{lemma}[Existence, uniqueness, and differentiability of thresholds]
\label{lem:generic_threshold_final}
Let $I\subset \mathbb R$ and $J=[\underline z,\bar z]\subset \mathbb R$ be
nonempty intervals, and let
\[
F:I\times J\to \mathbb R
\]
be continuous on $I\times J$ and continuously differentiable on its interior.
Suppose that
\begin{equation}
\frac{\partial F(u,z)}{\partial z}<0
\qquad
\text{for all }(u,z)\in I\times J.
\label{eq:Fz_negative_final}
\end{equation}
Then:

\begin{enumerate}
    \item For each fixed $u\in I$, the equation $F(u,z)=0$ has at most one
    solution in $J$.

    \item If, in addition,
    \begin{equation}
    F(u,\underline z)\ge 0
    \qquad\text{and}\qquad
    F(u,\bar z)\le 0,
    \label{eq:crossing_condition_final}
    \end{equation}
    then there exists a unique threshold $z^*(u)\in J$ such that
    \[
    F(u,z^*(u))=0.
    \]

    \item If the inequalities in \eqref{eq:crossing_condition_final} are
    strict, then $z^*(u)$ is continuously differentiable and
    \begin{equation}
    \frac{dz^*(u)}{du}
    =
    -\frac{\partial F(u,z^*(u))/\partial u}
    {\partial F(u,z^*(u))/\partial z}.
    \label{eq:implicit_formula_final}
    \end{equation}
\end{enumerate}
\end{lemma}

\begin{proof}
Condition \eqref{eq:Fz_negative_final} implies that, for each fixed $u$, the
map $z\mapsto F(u,z)$ is strictly decreasing, which gives uniqueness. If
\eqref{eq:crossing_condition_final} holds, continuity yields existence by the
intermediate value theorem. The differentiability statement follows from the
implicit function theorem because $\partial F/\partial z\neq 0$ at the
threshold.
\end{proof}

\subsection{Application to the threshold functions in Section~5}

We now apply Lemma~\ref{lem:generic_threshold_final} on the domain
$\mathcal D'$, where the monotonicity inequalities used in Section~5 are
numerically verified.

\noindent\textbf{Accession thresholds.}
Fix $\tau\in[0,\bar\tau]$. If
\begin{align}
\Delta_A^I(\tau,0) &\ge 0, &
\Delta_A^I(\tau,\bar\tau_C) &\le 0, \label{eq:AI_crossing_final}\\
\Delta_C^I(\tau,0) &\ge 0, &
\Delta_C^I(\tau,\bar\tau_C) &\le 0, \label{eq:CI_crossing_final}
\end{align}
then there exist unique thresholds $\tau_C^A(\tau),\tau_C^C(\tau)\in[0,\bar\tau_C]$
such that
\begin{align}
\Delta_A^I(\tau,\tau_C^A(\tau)) &= 0, \label{eq:tauCA_final}\\
\Delta_C^I(\tau,\tau_C^C(\tau)) &= 0. \label{eq:tauCC_final}
\end{align}
Moreover, if the endpoint inequalities in
\eqref{eq:AI_crossing_final}--\eqref{eq:CI_crossing_final} are strict, then
\begin{align}
\frac{d\tau_C^A(\tau)}{d\tau}
&=
-\frac{\partial \Delta_A^I(\tau,\tau_C^A(\tau))/\partial \tau}
{\partial \Delta_A^I(\tau,\tau_C^A(\tau))/\partial \tau_C},
\label{eq:d_tauCA_final}
\\
\frac{d\tau_C^C(\tau)}{d\tau}
&=
-\frac{\partial \Delta_C^I(\tau,\tau_C^C(\tau))/\partial \tau}
{\partial \Delta_C^I(\tau,\tau_C^C(\tau))/\partial \tau_C}.
\label{eq:d_tauCC_final}
\end{align}

\medskip
\noindent\textbf{Formation thresholds.}
Fix $\tau_C\in[0,\bar\tau_C]$. If
\begin{align}
\Delta_A^{SU}(0,\tau_C) &\ge 0, &
\Delta_A^{SU}(\bar\tau,\tau_C) &\le 0, \label{eq:ASU_crossing_final}\\
\Delta_A^{IS}(0,\tau_C) &\ge 0, &
\Delta_A^{IS}(\bar\tau,\tau_C) &\le 0, \label{eq:AIS_crossing_final}
\end{align}
then there exist unique thresholds $\tau^{SU}(\tau_C),\tau^{IS}(\tau_C)\in[0,\bar\tau]$
such that
\begin{align}
\Delta_A^{SU}(\tau^{SU}(\tau_C),\tau_C) &= 0, \label{eq:tauSU_final}\\
\Delta_A^{IS}(\tau^{IS}(\tau_C),\tau_C) &= 0. \label{eq:tauIS_final}
\end{align}
Moreover, if the endpoint inequalities in
\eqref{eq:ASU_crossing_final}--\eqref{eq:AIS_crossing_final} are strict, then
\begin{align}
\frac{d\tau^{SU}(\tau_C)}{d\tau_C}
&=
-\frac{\partial \Delta_A^{SU}(\tau^{SU}(\tau_C),\tau_C)/\partial \tau_C}
{\partial \Delta_A^{SU}(\tau^{SU}(\tau_C),\tau_C)/\partial \tau},
\label{eq:d_tauSU_final}
\\
\frac{d\tau^{IS}(\tau_C)}{d\tau_C}
&=
-\frac{\partial \Delta_A^{IS}(\tau^{IS}(\tau_C),\tau_C)/\partial \tau_C}
{\partial \Delta_A^{IS}(\tau^{IS}(\tau_C),\tau_C)/\partial \tau}.
\label{eq:d_tauIS_final}
\end{align}

\begin{corollary}[Threshold representation used in Section~5]
\label{cor:threshold_representation_final}
Under the monotonicity inequalities verified on $\mathcal D'$ and the crossing
conditions above, the primitive threshold functions
\[
\tau_C^A(\tau),\quad \tau_C^C(\tau),\quad \tau^{SU}(\tau_C),\quad \tau^{IS}(\tau_C)
\]
are well defined, unique, and continuously differentiable. Moreover,
\[
\tau_C^{I}(\tau)=\min\{\tau_C^A(\tau),\tau_C^C(\tau)\}
\]
and
\[
I_C^{I}(I_H)
=
\max\left\{
0,\,
\frac{\bar\tau_C-\tau_C^{I}(\bar\tau-\mu I_H)}{\mu}
\right\}
\]
are well defined and continuous, and piecewise continuously differentiable.
\end{corollary}

\medskip
\noindent\textbf{Scope of the threshold representation.}
Corollary~\ref{cor:threshold_representation_final} is the exact sense in which
threshold objects are used in Section~5. The paper does not claim that these
thresholds exist globally for all parameter values in $\mathcal P$. Rather,
the claim is that, on the verified compact domain $\mathcal D'$ and under the
endpoint crossing conditions stated above, the relevant regime-comparison
functions cross zero in a well-behaved way and therefore admit unique
differentiable threshold representations.

\subsection{Optional comparative-static refinement}

Using \eqref{eq:dAI_tau_v0_final}, if the domain $\mathcal D'$ is further
restricted so that
\[
\frac{\partial \Delta_A^I}{\partial \tau}<0
\qquad\text{on }\mathcal D',
\]
then \eqref{eq:d_tauCA_final} implies
\[
\frac{d\tau_C^A(\tau)}{d\tau}<0.
\]
Thus, a higher residual compliance cost for incumbent members lowers the
maximum outsider residual cost consistent with accession.
